\section{Results and Analysis}
\label{sec:results}
\draftnote{The independent rerun is complete. This section awaits integration
of the audited analysis. No historical comparison is used as a final result.
The red slots below must be replaced with the complete, version-matched
results and their qualified interpretation before submission.}

\subsection{Main Comparison}
% Replace with tools/update_results.py only after all three main tags complete.
\begin{table}[!t]
\caption{Main comparison result slots for ALNS and GCN-ALNS}
\label{tab:main_results}
\centering
\footnotesize
\begin{tabular}{llcc}
\toprule
Scale & Solver & Feasible (\%) & Cost, feasible \\
\midrule
Small & ALNS & \TBD & \TBD \\
      & GCN-ALNS & \TBD & \TBD \\
Medium & ALNS & \TBD & \TBD \\
       & GCN-ALNS & \TBD & \TBD \\
Large & ALNS & \TBD & \TBD \\
      & GCN-ALNS & \TBD & \TBD \\
\bottomrule
\end{tabular}
\end{table}

Table~\ref{tab:main_results} is reserved for the complete independent
comparison. The final presentation will include the entire baseline suite,
per-scale feasibility, and conditional cost. A companion paired table should
report GCN-ALNS minus ALNS on the shared evaluation population and on the
matched successful runs.

\pending{Insert the final per-scale comparison, confidence intervals, and
paired effect sizes. State explicitly whether GCN changes feasibility,
conditional cost, or recovery margin, and report any scale where the effect
reverses. Do not infer three-scale dominance from the completed small study.}

\subsection{Mechanism and Component Analysis}
The component results should separate changes to search behavior from
changes to the protection or recovery policy. In particular, the dynamic
rescheduling comparison measures the value of allowing cross-domain recovery
under the evaluator, whereas the feasibility-first and restart comparisons
measure search mechanisms. Removing all learned guidance measures the
increment supplied by the trained prior on top of those mechanisms.

\pending{Insert the 14-configuration component results and adjusted
significance values. Present domain-separation relaxation as a separate
admissibility or robustness comparison. Discuss learned subcomponents
without assuming that individually nonsignificant effects prove redundancy.}

\draftnote{A post-run audit found that the configurations labelled all-active
and all-passive do not enforce those modes for every task. Their current
results must not be presented as a strict pure-policy comparison. Describe
them as repair-mode preferences, or replace them with separately versioned
globally constrained experiments before making a hybrid-policy claim.}

\subsection{Iteration Budget and Runtime}
The budget sweep evaluates feasible rate and feasible cost against the
iteration limit. This experiment can establish whether learned guidance is
useful at restricted iteration budgets. Runtime must be read separately,
because learned construction and insertion can change per-iteration work and
the number of full-verification calls.

\pending{Insert the complete 50--1000 iteration curves and full solver
runtime comparison. Use ``first-feasible iteration'' for iteration counts;
do not rename an iteration-derived approximation as measured wall-clock
time-to-feasibility.}

